\spread{Writing a sentence with a run of visible glyphs}
\noindent The name of a battle can travel from a question to an answer in a single copy action. A question word can travel just as easily. We tested whether learning to copy longer runs helps a writer state the requested fact in a complete sentence.
\platehere{Figure_165_writing_with_visible_phrases}{A copy action can preserve a name or repeat the question}{Original answer pixels from the first held English case gained by the variable-run model and the first lost to the single-glyph model. Both saw the same pictured question and page. In panel a the run writer supplies the date. In panel b it copies the question through its final question mark and adds a period. Spaces at line breaks have been adjusted for print.}
\begin{multicols}{2}
\section*{What the action changes}
The writer can spell one bitmap, stop or copy 1, 2, 4, 8 or 16 contiguous cells from the question or a selected region. Origins producing identical runs share their probability mass. Each emitted glyph updates the recurrent state. We taught the longest available aligned run at every reference prefix and kept the same training order, region selector, initialization and four passes through the human replies. The extra length head adds 805 parameters. The copy mechanism operates directly on ink, including unfamiliar source glyphs. Span copying in token-sequence editing is an established antecedent (\href{https://arxiv.org/abs/2006.04771}{Panthaplackel et al.}).
\section*{Finding the value and completing the answer}
On 192 new English questions, fact presence rose from 44 to 62. Exact stopped statements fell from 14 to three. Both confidence intervals exclude zero when resampling the 96 passages (fact presence, +2.1 to +16.7 percentage points; exact images, $-$9.9 to $-$2.1). Both facts survived on 11 passages, up from four. Only the single-glyph writer recovered both complete English images on any passage, once. Disabling runs at inference in the new writer recovered ten exact English sentences. The longest-run writer also carried both versions of a changed fact on 22/62 English pairs, up from 17. In 35 new Chinese catalogue queries both writers gave nine complete human images. On 192 new Chinese examination questions they gave zero exact reference images.
\section*{Which training signal failed?}
For each stopped reference we summed the probabilities of all compatible ways to write the same ink. All 192 English references had greater probability than their supervised segmentation alone. Even after summing, mean loss per glyph rose from 0.319 to 0.417 nats relative to single-cell training. On 184 different stopped English replies, the run writer preferred the human reference to its own output only 18 times. A blind semantic judge accepted all 16 original English human controls but only five after the evidence was edited. The prespecified paired semantic score therefore remains unqualified.
\end{multicols}
\newpage
